\chapter{PHƯƠNG PHÁP ĐỀ XUẤT}

\section{Mô Hình}

\subsection{EfficientNet-B0 Backbone}

Kiến trúc backbone được chọn là EfficientNet-B0 \cite{tan2019efficientnet}, nạp từ thư viện timm \cite{wightman2019timm} với trọng số tiền huấn luyện trên ImageNet:

\begin{lstlisting}[style=python, caption={Khởi tạo EfficientNet-B0 backbone}]
import timm
backbone = timm.create_model(
    "efficientnet_b0",
    pretrained=True,
    num_classes=0,   # loai bo classifier mac dinh
    global_pool=""   # them pooling rieng
)
# output: feature map (B, 1280, 7, 7)
\end{lstlisting}

EfficientNet được xây dựng trên nguyên tắc \textit{compound scaling} --- scale đồng thời độ sâu (depth), độ rộng (width) và độ phân giải (resolution) theo một hệ số tỷ lệ thống nhất. EfficientNet-B0 --- phiên bản nhỏ nhất --- đạt accuracy tương đương ResNet-50 với $\approx 5{,}3$M tham số (so với $\approx 25{,}6$M), phù hợp fine-tuning trên tập dữ liệu đa domain.

\begin{table}[H]
  \centering
  \caption{So sánh kiến trúc backbone}
  \label{tab:backbone_compare}
  \begin{tabular}{llrl}
    \toprule
    \textbf{Kiến trúc} & \textbf{Tham số} & \textbf{Top-1 Acc. (ImageNet)} & \textbf{Ghi chú} \\
    \midrule
    EfficientNet-B0 & $\approx$5,3M & 77,1\% & \textit{Lựa chọn đề tài} \\
    ResNet-50       & $\approx$25,6M & 76,2\% & $\approx 5\times$ tham số hơn \\
    VGG-16          & $\approx$138M  & 74,4\% & Quá nặng cho fine-tuning \\
    ViT-B/16        & $\approx$86M   & 81,8\% & Cần nhiều dữ liệu hơn \\
    \bottomrule
  \end{tabular}
\end{table}

\subsection{Custom Classification Head}

Sau backbone, một phần đầu phân loại tùy chỉnh được gắn vào:

\begin{figure}[H]
  \centering
  \begin{minipage}{0.55\linewidth}
    \begin{lstlisting}[style=python, caption={Kiến trúc classification head}]
nn.Sequential(
    nn.AdaptiveAvgPool2d(1),  # (B,1280,1,1)
    nn.Flatten(),              # (B, 1280)
    nn.Linear(1280, 256),
    nn.ReLU(inplace=True),
    nn.Dropout(p=0.5),
    nn.Linear(256, 1)          # logit
)
    \end{lstlisting}
  \end{minipage}
\end{figure}

\begin{table}[H]
  \centering
  \caption{Thiết kế và lý do từng thành phần trong classification head}
  \label{tab:head}
  \begin{tabularx}{\linewidth}{llX}
    \toprule
    \textbf{Thành phần} & \textbf{Cấu hình} & \textbf{Lý do thiết kế} \\
    \midrule
    AdaptiveAvgPool2d & Output (1,1) & Tổng hợp không gian từ toàn bộ feature map \\
    Linear (đầu tiên) & $1280 \to 256$ & Đủ rộng để học ranh giới phân loại Real/Fake \\
    Dropout           & $p = 0{,}5$   & Ngăn overfitting --- artifact dễ bị ``học thuộc'' \\
    Linear (đầu ra)   & $256 \to 1$   & Một logit cho phân loại nhị phân \\
    \bottomrule
  \end{tabularx}
\end{table}

\textbf{Hàm mất mát:} \texttt{BCEWithLogitsLoss} kết hợp sigmoid và binary cross-entropy:
\begin{equation}
  \mathcal{L}(x, y) = -\bigl[ y \cdot \log\sigma(x) + (1-y) \cdot \log(1 - \sigma(x)) \bigr]
  \label{eq:bce}
\end{equation}
trong đó $x$ là logit đầu ra, $y \in \{0, 1\}$ là nhãn thật, $\sigma(\cdot)$ là hàm sigmoid. Sigmoid \textbf{không} được áp dụng trước khi tính mất mát --- chỉ áp dụng khi inference.

\subsection{Sơ Đồ Kiến Trúc Tổng Thể}

Hình~\ref{fig:model_architecture} trực quan hóa luồng dữ liệu qua toàn bộ mô hình từ ảnh đầu vào đến quyết định phân loại cuối cùng.

\begin{figure}[H]
  \centering
  \begin{tikzpicture}[
    node distance=0.52cm,
    every node/.style={font=\small},
    bk/.style  = {rectangle, draw=orange!70!black, fill=orange!7, rounded corners=3pt,
                  minimum width=9.5cm, minimum height=0.88cm, align=center},
    hd/.style  = {rectangle, draw=green!60!black, fill=green!7, rounded corners=3pt,
                  minimum width=9.5cm, minimum height=0.88cm, align=center},
    io/.style  = {rectangle, draw=uitblue, fill=uitblue!8, rounded corners=3pt,
                  minimum width=9.5cm, minimum height=0.88cm, align=center},
    arr/.style = {-Stealth, semithick}
  ]
    \node[io] (inp)  {Ảnh đầu vào $\;\big|\;$ Resize($224{\times}224$) $\;\big|\;$ Normalize (ImageNet)};
    \node[bk, below=of inp]  (bk0)
      {EfficientNet-B0: MBConv Blocks 0--4 \quad \textit{(đóng băng ở Giai đoạn 1--2)}};
    \node[bk, below=of bk0]  (bk1)
      {EfficientNet-B0: MBConv Blocks 5--6 \quad \textit{(mở băng từ Giai đoạn 2)}};
    \node[bk, below=of bk1]  (fm)
      {Feature Map $\quad 1280 \times 7 \times 7$};
    \node[hd, below=of fm]   (gp)
      {AdaptiveAvgPool2d(1,1) $\;\to\;$ Flatten $\;\to\;$ $\mathbb{R}^{1280}$};
    \node[hd, below=of gp]   (fc1)
      {Linear(1280, 256) $\;\to\;$ ReLU $\;\to\;$ Dropout(0.5)};
    \node[hd, below=of fc1]  (fc2)
      {Linear(256, 1) \quad Logit $\in \mathbb{R}$};
    \node[io, below=of fc2]  (out)
      {Sigmoid $\to$ $p \in [0,1]$ $\;\;\big|\;\;$ Real ($p \leq 0{,}5$) $\;/\;$ Fake ($p > 0{,}5$)};

    \foreach \a/\b in {inp/bk0, bk0/bk1, bk1/fm, fm/gp, gp/fc1, fc1/fc2, fc2/out}
      \draw[arr] (\a.south) -- (\b.north);
  \end{tikzpicture}
  \caption{Kiến trúc tổng thể mô hình EfficientNet-B0 với classification head tùy chỉnh.
           Cam: EfficientNet-B0 backbone pretrained. Xanh lá: custom classification head.}
  \label{fig:model_architecture}
\end{figure}

\subsection{Chiến Lược Huấn Luyện 3 Giai Đoạn}

Chiến lược dựa trên nguyên tắc \textit{discriminative fine-tuning}: các lớp gần đầu vào đã học đặc trưng tốt từ ImageNet nên cần learning rate thấp; các lớp đầu ra mới cần học từ đầu với learning rate cao hơn.

\begin{table}[H]
  \centering
  \caption{Thông số chi tiết 3 giai đoạn huấn luyện}
  \label{tab:stages}
  \begin{tabularx}{\linewidth}{clXrrl}
    \toprule
    \textbf{GĐ} & \textbf{Tên} & \textbf{Lớp trainable} & \textbf{LR} & \textbf{Max E.} & \textbf{Checkpoint} \\
    \midrule
    1 & Freeze Backbone   & Head only                   & $10^{-4}$ & 30 & \texttt{best\_stage1.pth} \\
    2 & Unfreeze Top      & Head + blocks 5--6          & $10^{-5}$ & 30 & \texttt{best\_stage2.pth} \\
    3 & Unfreeze All      & Toàn bộ mô hình             & $10^{-6}$ & 30 & \texttt{best\_stage3.pth} \\
    \bottomrule
  \end{tabularx}
\end{table}

\begin{table}[H]
  \centering
  \caption{Kết quả huấn luyện trên tập kiểm định theo giai đoạn}
  \label{tab:train_results}
  \begin{tabular}{clrrrl}
    \toprule
    \textbf{GĐ} & \textbf{Epochs chạy} & \textbf{Val Loss} & \textbf{Val Accuracy} & \textbf{Kích thước} \\
    \midrule
    1 & 30/30 & 0,3103 & 86,63\% & 20,3 MB \\
    2 & 30/30 & 0,1383 & 94,54\% & 26,0 MB \\
    3 & 30/30 & 0,0699 & 97,34\% & 52,5 MB \\
    \bottomrule
  \end{tabular}
\end{table}

Đáng chú ý, EarlyStopping không kích hoạt ở bất kỳ giai đoạn nào --- mô hình tiếp tục cải thiện đến epoch cuối. Val loss giảm liên tục $0{,}310 \to 0{,}138 \to 0{,}070$ ($\approx 77\%$ tổng mức giảm).

Hình~\ref{fig:training_stages} trình bày sơ đồ luồng toàn bộ quy trình huấn luyện 3 giai đoạn, từ khởi tạo mô hình đến đánh giá trên tập test độc lập.

\begin{figure}[H]
  \centering
  \begin{tikzpicture}[
    every node/.style={font=\small},
    stage/.style = {rectangle, draw=uitblue, fill=uitblue!10, rounded corners=4pt,
                    minimum width=3.4cm, minimum height=2.6cm, align=center},
    res/.style   = {rectangle, draw=green!60!black, fill=green!7, rounded corners=3pt,
                    text width=3.2cm, minimum height=0.88cm, align=center, font=\footnotesize},
    fin/.style   = {rectangle, draw=orange!70!black, fill=orange!8, rounded corners=4pt,
                    minimum width=11.5cm, minimum height=0.95cm, align=center},
    arr/.style   = {-Stealth, semithick, uitblue}
  ]
    \node[stage] (s1) at (0,0) {
      \textbf{Giai đoạn 1}\\[4pt]
      \small Freeze Backbone\\
      LR $= 10^{-4}$\\
      30 epochs\\
      Head only
    };
    \node[stage] (s2) at (4.2,0) {
      \textbf{Giai đoạn 2}\\[4pt]
      \small Unfreeze B5--B6\\
      LR $= 10^{-5}$\\
      30 epochs\\
      Head + Blocks 5--6
    };
    \node[stage] (s3) at (8.4,0) {
      \textbf{Giai đoạn 3}\\[4pt]
      \small Full Fine-tune\\
      LR $= 10^{-6}$\\
      30 epochs\\
      Toàn bộ mô hình
    };

    \node[res, below=0.5cm of s1] (r1) {Val Acc: 86,63\%\\Val Loss: 0,310};
    \node[res, below=0.5cm of s2] (r2) {Val Acc: 94,54\%\\Val Loss: 0,138};
    \node[res, below=0.5cm of s3] (r3) {Val Acc: 97,34\%\\Val Loss: 0,070};

    % eval positioned directly below middle result — no empty space
    \node[fin, below=0.75cm of r2] (eval) {
      Đánh giá trên Test Set (47.480 ảnh):
      \quad Acc 97,33\% $\;|\;$ F1 97,45\% $\;|\;$ AUC-ROC 99,70\%
    };

    % Horizontal arrows — label slightly raised to clear box border
    \draw[arr] (s1.east) -- node[above, font=\footnotesize\itshape, yshift=1pt] {checkpoint} (s2.west);
    \draw[arr] (s2.east) -- node[above, font=\footnotesize\itshape, yshift=1pt] {checkpoint} (s3.west);

    % Vertical arrows from stages to results
    \foreach \s/\r in {s1/r1, s2/r2, s3/r3}
      \draw[-Stealth, semithick] (\s.south) -- (\r.north);

    % Three vertical arrows converging straight down into eval.north
    \draw[-Stealth, semithick] (r1.south) -- (r1.south |- eval.north);
    \draw[-Stealth, semithick] (r2.south) -- (eval.north);
    \draw[-Stealth, semithick] (r3.south) -- (r3.south |- eval.north);
  \end{tikzpicture}
  \caption{Sơ đồ chiến lược huấn luyện 3 giai đoạn progressive unfreezing và kết quả từng giai đoạn}
  \label{fig:training_stages}
\end{figure}

\section{Pipeline Tiền Xử Lý và Tăng Cường Dữ Liệu}

\subsection{Tiền Xử Lý Cơ Bản}

Mọi ảnh đầu vào trải qua hai bước cố định:
\begin{enumerate}
  \item \textbf{Resize về $224 \times 224$:} Kích thước đầu vào mặc định của EfficientNet-B0. EDA xác nhận hầu hết ảnh đã ở tỷ lệ khung hình $\approx 1{:}1$, không gây biến dạng.
  \item \textbf{Chuẩn hóa ImageNet:} Mean $= [0{,}485;\ 0{,}456;\ 0{,}406]$, Std $= [0{,}229;\ 0{,}224;\ 0{,}225]$ (kênh R/G/B).
\end{enumerate}

\subsection{Tăng Cường Dữ Liệu (Train Only)}

\begin{table}[H]
  \centering
  \caption{Các phép biến đổi tăng cường dữ liệu (albumentations)}
  \label{tab:augmentation}
  \begin{tabularx}{\linewidth}{lXrl}
    \toprule
    \textbf{Phép biến đổi} & \textbf{Tham số} & \textbf{$p$} & \textbf{Mục đích} \\
    \midrule
    HorizontalFlip & --- & 0,5 & Mô phỏng đối xứng gương \\
    RandomBrightnessContrast & brightness\_limit=0,2; contrast\_limit=0,2 & 0,5 & Mô phỏng điều kiện ánh sáng \\
    ShiftScaleRotate & shift=0,1; scale=0,1; rotate=15° & 0,5 & Mô phỏng góc chụp thay đổi \\
    CoarseDropout & max\_holes=8; kích thước $16{\times}16$~px & 0,3 & Ngăn phụ thuộc đặc trưng cục bộ \\
    \bottomrule
  \end{tabularx}
\end{table}

Tập val/test chỉ áp dụng resize và chuẩn hóa --- không augmentation để đảm bảo đánh giá tái lập.

\section{Môi Trường Huấn Luyện và Công Nghệ Sử Dụng}

Mô hình được huấn luyện trong $\approx 17$ giờ trên phần cứng dưới đây.

\begin{table}[H]
  \centering
  \caption{Cấu hình phần cứng huấn luyện}
  \label{tab:hardware}
  \begin{tabular}{ll}
    \toprule
    \textbf{Thành phần} & \textbf{Thông số} \\
    \midrule
    CPU  & AMD Ryzen 5 7500F 6-Core \\
    GPU  & NVIDIA GeForce RTX 5070 \\
    VRAM & 12,82 GB \\
    Compute Capability & 12.0 (Blackwell) \\
    CUDA & 13.0 \\
    \bottomrule
  \end{tabular}
\end{table}

\begin{table}[H]
  \centering
  \caption{Các thư viện và phiên bản sử dụng}
  \label{tab:libraries}
  \begin{tabular}{llll}
    \toprule
    \textbf{Loại} & \textbf{Thư viện} & \textbf{Phiên bản} & \textbf{Mục đích} \\
    \midrule
    Framework     & PyTorch       & 2.12.0+cu130 & Deep learning \\
    Model library & timm          & 1.0.27       & EfficientNet-B0 pretrained \\
    Augmentation  & albumentations & 2.0.8       & Train transform \\
    Metrics       & scikit-learn  & 1.8.0        & Acc/P/R/F1/AUC-ROC \\
    Visualization & matplotlib    & $\geq$3.8    & Confusion matrix, ROC curve \\
    Data          & pandas        & $\geq$2.1    & Đọc và xử lý CSV \\
    Config        & PyYAML        & ---          & Nạp siêu tham số từ YAML \\
    Ngôn ngữ      & Python        & 3.12         & --- \\
    \bottomrule
  \end{tabular}
\end{table}

\textbf{Mixed-precision training (FP16):} Được bật bằng \texttt{torch.cuda.amp.autocast()} kết hợp \texttt{GradScaler}, giảm VRAM $\approx 40\%$ và tăng tốc $\approx 1{,}5$--$2\times$ trên RTX 5070.

\textbf{Tái lập kết quả:} Seed 42 cố định cho \texttt{random}, \texttt{numpy}, \texttt{torch} và \texttt{torch.cuda}; \texttt{deterministic=True}.

\section{Lớp FaceDataset và DataLoader}

\subsection{Thiết Kế Lớp FaceDataset}

Lớp \texttt{FaceDataset} (định nghĩa tại \texttt{src/dataset.py}) kế thừa \texttt{torch.utils.data.Dataset} và đóng vai trò giao diện giữa file CSV phân chia và pipeline huấn luyện. Luồng xử lý bên trong \texttt{\_\_getitem\_\_} gồm ba bước:

\begin{enumerate}
  \item Đọc đường dẫn ảnh và nhãn từ DataFrame đã nạp trong \texttt{\_\_init\_\_}.
  \item Mở ảnh bằng \texttt{PIL.Image.open()}, chuyển sang mảng NumPy RGB, áp dụng transform (albumentations hoặc chỉ resize+normalize cho val/test).
  \item Trả về cặp \texttt{(image\_tensor, label\_tensor)} với \texttt{label\_tensor} kiểu \texttt{torch.float32} --- kiểu dữ liệu bắt buộc bởi \texttt{BCEWithLogitsLoss}.
\end{enumerate}

\subsection{Cấu Hình DataLoader}

\begin{table}[H]
  \centering
  \caption{Cấu hình DataLoader cho ba tập dữ liệu}
  \label{tab:dataloader}
  \begin{tabular}{llll}
    \toprule
    \textbf{Tham số} & \textbf{Train} & \textbf{Val} & \textbf{Test} \\
    \midrule
    batch\_size   & 128  & 128  & 128  \\
    shuffle       & True & False & False \\
    num\_workers  & 2    & 2    & 2    \\
    pin\_memory   & True & True & True \\
    drop\_last    & False & False & False \\
    \bottomrule
  \end{tabular}
\end{table}

\textbf{Lý do chọn \texttt{num\_workers=2}:} Trên hệ thống với CPU 6 nhân, \texttt{num\_workers=4} gây tranh chấp tài nguyên với tiến trình huấn luyện chính. Qua thử nghiệm, giá trị 2 cho thời gian nạp dữ liệu đủ nhanh để GPU không bị đợi (bottleneck tại data loading $<$ 5\% thời gian mỗi batch).

\subsection{Quy Trình Quản Lý Siêu Tham Số}

Toàn bộ siêu tham số được khai báo tập trung trong \texttt{configs/train\_config.yaml}. Không có giá trị nào được hardcode trong mã nguồn. Khi khởi chạy, script đọc file YAML, ghi log đầy đủ các tham số cùng với phiên bản thư viện và thông tin GPU, đảm bảo mọi lần chạy đều có thể tái lập và truy xuất cấu hình sau này.

Quy trình quản lý thực nghiệm này --- seed cố định, tham số qua YAML, log phiên bản --- là thực hành tốt trong nghiên cứu học sâu để đảm bảo tính tái lập (reproducibility), một yêu cầu ngày càng được nhấn mạnh trong cộng đồng ML nghiên cứu.
